\chapter{Background and Literature}
\label{ch:background_literature}

\section{Healthcare Access Theory}

Unmet medical care need is an access indicator, not simply a utilisation indicator. The classical access literature treats realised care as the outcome of demand, resources, health-system organisation, and perceived need. Andersen's behavioural model and the Aday-Andersen framework distinguish predisposing characteristics, enabling resources, need, and health-system organisation as determinants of realised care \citep{aday_andersen_1974_framework, andersen_1995_behavioral_model}. Penchansky and Thomas define access through availability, accessibility, affordability, accommodation, and acceptability \citep{penchansky_thomas_1981_access}. Levesque, Harris, and Russell extend this view by modelling access as an interface between health-system characteristics and people's abilities to perceive, seek, reach, pay for, and engage with care \citep{levesque_harris_russell_2013_access}.

Those frameworks imply that unmet care is not produced by one mechanism. A person can report unmet need because care is unaffordable, because a service is geographically or organisationally inaccessible, because waiting times are long, because administrative navigation is difficult, because expectations differ, or because need itself is perceived differently. This thesis cannot observe those mechanisms at the individual level. It uses country-year aggregates, so it translates access theory into broad monitoring domains rather than into individual causal pathways.

The variable mapping follows that restricted purpose. Poverty or social exclusion and out-of-pocket spending are closest to affordability. Physicians and hospital beds are coarse availability proxies. GDP per capita captures macroeconomic resources. Unemployment captures labour-market vulnerability. Country effects stand in for persistent institutional arrangements, while year effects capture common period shocks. This mapping justifies why these variables belong in an aggregate monitoring model; it does not make the fixed-effects coefficient causal.

\begin{figure}[htbp]
\centering
\includegraphics[width=0.92\textwidth]{conceptual_selection_estimand_diagram}
\caption{Conceptual selection and estimand diagram for the aggregate Eurostat design. The figure is not a causal DAG. It shows how access-domain covariates, country/year structure, missing covariate availability, and complete-case inclusion jointly determine which monitoring estimand is being estimated.}
\label{fig:conceptual_selection_estimand_diagram}
\end{figure}

Figure~\ref{fig:conceptual_selection_estimand_diagram} gives the chapter's organising logic. Social vulnerability, macroeconomic context, and health-system proxies are related to the observed aggregate outcome, but missing covariate availability also determines whether a country-year enters the complete-case regression. Missingness is therefore part of the estimand, not only a nuisance.

\section{Measurement of Self-Reported Unmet Care}

Eurostat's EU-SILC health metadata state that the health component contains health-status variables and unmet-needs variables for medical and dental care \citep{eurostat_health_variables_metadata}. The unmet-needs variables refer to a respondent's own assessment of whether they needed a specific type of examination or treatment but did not have it or did not seek it \citep{eurostat_health_variables_silc_methodology}. Medical care excludes dental care and refers to medical examination or treatment provided by, or under the supervision of, medical doctors or equivalent professions \citep{eurostat_health_variables_silc_methodology}.

The source is harmonised but not mechanically identical to an administrative registry. Eurostat describes EU-SILC methodology as being defined through a legislative framework to support comparability across EU, EFTA, and candidate countries, with yearly methodological guidelines for implementation \citep{eurostat_eusilc_methodology}. At the same time, EU-SILC combines interview data, national survey implementation, administrative sources, and imputations in ways that require explicit documentation \citep{eurostat_silc_information_data}. The empirical design therefore treats the Eurostat table as a public monitoring source whose comparability is supported institutionally, while still requiring checks for coverage, denominator choice, and measurement-period changes.

The official health-variable methodology lists the main reasons for unmet need as categories including care being too expensive, waiting list, inability to take time because of work or care responsibilities, too far to travel or no transportation, fear, waiting to see whether the problem improved, not knowing a good doctor or dentist, and other reasons \citep{eurostat_health_variables_silc_methodology}. This thesis uses the Eurostat aggregate category that combines the access-barrier reasons most directly connected to system functioning: expense, distance, and waiting lists. These are substantively related, but they are not identical mechanisms.

The primary outcome, \texttt{hlth\_silc\_08}, is an all-population percentage of people aged 16 or over reporting unmet medical examination or treatment. Eurostat's unmet-healthcare article distinguishes this type of indicator from a second indicator expressed as a percentage of people who reported having a medical need \citep{eurostat_unmet_needs_statistics}. The robustness outcome, \texttt{hlth\_silc\_08b}, belongs to that second denominator family. It is therefore useful as a denominator check, but it is not an interchangeable replacement for \texttt{hlth\_silc\_08}.

This denominator distinction is analytically important. A population-denominator rate combines at least two components: whether people perceive or report a need and whether that need is unmet. A need-denominator rate focuses more narrowly on unmet access among people reporting need. \citet{A4_ramos_2019} argue that ordinary unmet-need measures can hide this difference, and \citet{ingleby_guidi_2025_unmet_needs_methodology} explain that the 2015 EU-SILC change made it possible to distinguish the probability of reporting need from the probability that care was not received. This is why the thesis treats \texttt{hlth\_silc\_08b} as measurement robustness over 2021--2025 rather than as a second main outcome.

EU-SILC is also a survey of private households and persons living in those households. Eurostat states that all household members are surveyed but only those aged 16 or older are interviewed, while people in collective households and institutions are generally outside the target population \citep{eurostat_silc_information_data, eurostat_income_living_metadata}. Health variables are collected through interviews rather than administrative records \citep{eurostat_health_variables_silc_methodology}. Consequently, self-reported unmet care should not be read as clinical need, administrative denial of service, or objective health-system quality. It is a harmonised monitoring indicator based on reported need, reported non-use, and reason categories.

The 2015 issue is handled as a measurement-period caution. The literature notes that a change in the EU-SILC health-needs questionnaire from 2015 onwards separated the need question from the unmet-care question, making need-denominator indicators possible \citep{ingleby_guidi_2025_unmet_needs_methodology}. In the present aggregate panel, the baseline complete-case regression starts in 2015 because covariate coverage is thin before that year. A formal complete-case pre/post-2015 split is therefore not available for the main poverty coefficient. Chapter~\ref{ch:modeling_strategy_results} uses excluding-2015 and 2016-onward checks, while descriptive pre/post patterns are interpreted cautiously.

\section{Individual-Level European Evidence}

Individual-level European studies consistently link unmet medical care need to social vulnerability. Lower income, unemployment, difficulty making ends meet, poor health, debt, and housing vulnerability are associated with higher access barriers in EU-SILC and related European survey evidence \citep{A1_chaupain_guillot_2015, A2_madureira_lima_2018, A3_israel_2016, A5_cylus_papanicolas_2015}. Studies of financing and protection show that out-of-pocket payments, social-policy design, and public pensions can matter for cost-related unmet need \citep{B2_kaminska_wulfgramm_2019, C4_reeves_mckee_mackenbach_2017}. Crisis-period work from Greece, Portugal, and the Baltic states links economic stress and health-system pressure to access problems \citep{zavras_zavras_kyriopoulos_2016, legido_quigley_karanikolos_2016, karanikolos_gordeev_mackenbach_2016}.

This evidence motivates the choice of poverty or social exclusion as the stress-test coefficient. It does not imply that the thesis estimates the individual-level social gradient. Microdata studies can observe household income, self-reported health, employment, education, citizenship, and personal characteristics. This thesis observes national averages. Agreement with microdata literature can make the aggregate poverty association substantively plausible, but it cannot convert a country-year coefficient into evidence about individual behaviour.

The same distinction applies to citizenship. European access studies often treat migration status, legal entitlement, language, and socioeconomic position as linked access dimensions. The aggregate Eurostat citizenship rows used in this thesis cannot adjust for income, health status, legal status, language, local service conditions, or duration of residence. They are therefore descriptive equity context, not a parallel contribution to individual-level access theory.

\section{Aggregate Monitoring and Ecological Inference}

Country-year panels are useful for monitoring because they make cross-national coverage, trends, and policy-relevant indicators visible. They are weak for mechanisms because the unit of analysis is the country-year rather than the person. Classic work on ecological correlations shows that aggregate relationships can differ from individual relationships \citep{robinson_1950}. Epidemiological work further shows that ecological bias can arise through aggregation, contextual confounding, compositional differences, and effect modification \citep{greenland_morgenstern_1989}.

This is not a minor caveat. Poverty, unemployment, citizenship, and unmet care operate through individual and household mechanisms, but this thesis observes only country-year summaries. A positive country-year poverty coefficient does not show that poorer individuals in a country forgo care. It may reflect compositional change, national policy context, reporting expectations, macroeconomic shocks, health-system pressure, or measurement coverage. For this reason, the estimand is explicitly an aggregate monitoring estimand.

The aggregate design is still defensible when the question is framed correctly. The thesis does not ask whether poverty causes an individual to report unmet need. It asks whether an aggregate Eurostat poverty association is stable when the analyst changes target population, weighting, row selection, and imputation assumptions. That is a methodological monitoring question rather than a micro-level access question.

\section{Missing Data, Complete Cases, and Estimand Drift}

Missingness is the central methodological problem because the outcome is observed for more country-years than the baseline covariates. Complete-case analysis can therefore change the target population. Missing-data theory distinguishes complete-case analysis, inverse-probability weighting, and multiple imputation under explicit assumptions \citep{rubin_1987, little_rubin_2019}. IPW can address missingness when observation probabilities are well modelled and positivity holds, but it cannot recover structurally absent covariates \citep{seaman_white_2013}.

In this thesis, estimand drift occurs when the analysis silently moves from the 608-row outcome-observed panel to the 282-row complete-case fixed-effects sample. That movement is not only a loss of power. It changes which countries and years define the coefficient. If excluded rows have different unmet need, GDP, unemployment, or poverty levels, then the complete-case coefficient belongs to a selected analytical panel rather than to the intended monitoring target.

Multiple imputation is used as a sensitivity design, not as a rescue operation. A MAR-type imputation model asks what happens if missing covariates can be predicted from observed variables, country or country-group structure, year structure, and the observed outcome. That assumption is not verifiable from the observed aggregate data. Structural Eurostat non-coverage, country-specific reporting practices, and variables that exist only over shorter windows may violate MAR. The MNAR delta-shift checks therefore ask whether the conclusion depends mechanically on one MAR imputation model.

The interpretation rule is deliberately conservative. If complete-case specifications remain positive but MAR and MNAR full-target imputation attenuate the coefficient, the thesis does not declare the imputed coefficient to be true. It concludes that the complete-case poverty association is estimand-dependent and model-dependent. This is the core contribution.

\section{Reporting Discipline for Observational Monitoring}

The thesis follows a STROBE-style reporting logic for observational data: data source, eligibility, variables, bias, statistical methods, missing data, and limitations are stated before interpreting model output \citep{von_elm_altman_egger_2007_strobe}. This matters because the project is not a controlled experiment and does not have a single natural denominator. The key reporting unit is therefore not only the regression coefficient, but also the target population, row-selection rule, country universe, weighting rule, covariance estimator, and missing-data assumption attached to each coefficient.

This reporting discipline is also why the thesis avoids treating robustness checks as interchangeable. Complete-case, population-weighted, balanced-outcome, IPW, MAR MI, and MNAR rows are different estimands or sensitivity targets. They do not form a simple tournament for the one correct answer. Their value is to reveal where the interpretation changes.

\section{Panel Dependence and Prediction Validation}

Panel inference requires dependence-aware uncertainty because each country contributes repeated observations. Cluster-robust inference for within estimators is standard when errors may be correlated within countries \citep{arellano_1987}. The panel literature also warns that serial correlation can lead to overconfident inference if repeated observations are treated as independent \citep{bertrand_duflo_mullainathan_2004}. For this reason, the thesis reports country-clustered standard errors for pooled and two-way fixed-effects models, with Driscoll--Kraay and wild-cluster bootstrap diagnostics added for the main selected-panel coefficient.

Prediction is secondary but must still respect the panel structure. Country-year data have temporal dependence and strong country persistence, so random row splits overstate generalisation. Structured validation is recommended when observations are temporally, spatially, or hierarchically dependent \citep{roberts_2017_cross_validation}. Recent prediction-model reporting and bias-assessment guidance reinforces the same principle: the target population, outcome, predictors, validation design, performance uncertainty, and risk of bias must be transparent \citep{collins_moons_dhiman_2024_tripod_ai, moons_damen_kaul_2025_probast_ai}. In this thesis those standards are used only to discipline the benchmark exercise, not to claim a clinical prediction-model contribution.

The thesis therefore defines prediction as one-year-ahead benchmarking using lagged predictors. Same-year prediction is closer to nowcasting and is not treated as the main predictive claim. Flexible machine-learning models are benchmark competitors, not a separate scientific contribution. Feature importance is model inspection only, because tree-model impurity importance can be biased by scale, variation, and split opportunities \citep{strobl_boulesteix_zeileis_hothorn_2007_bias}.
